%Это~"--- заготовка для перевода статьи о~User Journeys и~User Flows. https://www.nngroup.com/articles/user"=journeys"=vs"=user"=flows/ В~ходе перевода улетела в~кроличью нору перекрестных ссылок о~методологиях дизайн"=мышления. Сначала переведу Design Thinking.

\documentclass[a4paper,12pt]{article}
\usepackage[T2A]{fontenc}
\usepackage[english, main=russian]{babel}
\usepackage{pscyr}
\usepackage{hyperref}
\usepackage{graphicx}
\renewcommand{\familydefault}{fac}
\begin{document}

\title{User~Journeys vs. User~Flows}
\author{Kate Kaplan}
\date{актуально на 27~августа 2026}

\maketitle

\textit{User journeys} и~\textit{user flows}~"--- инструменты UX, которые охватывают то, как люди достигают целей с~помощью конкретных продуктов и~сервисов. У~них есть схожие черты. И~user journey и~user flow:

\begin{enumerate}
	\item Используются во время этапа
		\footnote{
			\textit{Ideation} или «генерация идей»"--- один из этапов применения дизайн"=мышления по методологии \textit{Hasso Plattner Institute of Design (d.school)}. Переводчица копала до первоисточников, а~оказалось, что у~NNG уже есть \href{https://www.nngroup.com/articles/design"=thinking/}{статья о~дизайн"=мышлении.} Упс! :) 
			\par
			Традиционно выделяют пять этапов дизайн"=мышления, NNG добавляет шестой: 
				\begin{enumerate}
					\item Empathize~"--- эмпатия
					\item Define~"--- анализ и~определение проблемы
					\item Ideate~"--- генерация идей
					\item Prototype~"--- прототипирование
					\item Test~"--- тестирование
					\item Implement~"--- внедрение
				\end{enumerate} 
			\href{https://designthinking.ideo.com/process}{IDEO выделяет семь других этапов}: \textit{Frame the question, Gather inspiration, Synthesize for action, Generate ideas, Make ideas tangible, Test to learn, Share the story.}
			\par
			Как большая поклонница мюзикла \textit{Wicked} и~теоретик, переводчица считает, что нужно упомянуть термин \textit{\href{https://ixdf.org/literature/topics/wicked"=problems}{Wicked Problems}}. Это термин 1960"=x,~описывающий сложные задачи, которые сложно определить, у~которых нет единого решения, явного финала и~фиксированных требований. Хорст Риттель и~Мелвин Уэббер сформулировали десять принципов \textit{\href{https://www.sympoetic.net/Managing\_Complexity/complexity\_files/1973\%20Rittel\%20and\%20Webber\%20Wicked\%20Problems.pdf}{wicked problems}}, применимо в~контексте социального планирования. \href{https://www.researchgate.net/publication/251880127\_Wicked\_Problems\_in\_Design\_Thinking}{Ричард Буканан развил эту идею}, и~предположил, что к~ним относится большинство задач дизайнера. Это происходит, потому что каждый раз дизайнер переизобретает предмет дизайна для конкретной ситуации.
		}
		генерации идей или оценки, с~целью понять и~оптимизировать пользовательский опыт.
	\item Построены вокруг цели пользователя и~изучаются с~точки зрения пользователя или потребителя (а~не компании или продукта).
	\item Описываются и~визуализируются с~помощью карт.

\end{enumerate}

\end{document}